\documentclass[11pt,fontset=fandol]{ctexart}
\usepackage[a4paper,margin=23mm]{geometry}
\usepackage{amsmath,amssymb,mathtools,bm,mathrsfs,esint}
\usepackage{graphicx,longtable,booktabs,array,calc}
\usepackage[normalem]{ulem}
\usepackage{xcolor,hyperref}
\hypersetup{hidelinks}
\setcounter{secnumdepth}{-1}
\newcounter{none}
\providecommand{\tightlist}{\setlength{\itemsep}{0pt}\setlength{\parskip}{0pt}}
\providecommand{\passthrough}[1]{#1}
\usepackage{amssymb,mathrsfs,esint,bm,newunicodechar}
\allowdisplaybreaks
\setlength{\emergencystretch}{3em}
\providecommand{\llangle}{\langle\!\langle}
\providecommand{\rrangle}{\rangle\!\rangle}
\newunicodechar{∞}{\ensuremath{\infty}}
\newunicodechar{∈}{\ensuremath{\in}}
\newunicodechar{∑}{\ensuremath{\sum}}
\newunicodechar{∫}{\ensuremath{\int}}
\newunicodechar{≤}{\ensuremath{\leq}}
\newunicodechar{≥}{\ensuremath{\geq}}
\newunicodechar{≠}{\ensuremath{\neq}}
\newunicodechar{→}{\ensuremath{\to}}
\newunicodechar{⇒}{\ensuremath{\Rightarrow}}
\newunicodechar{α}{\ensuremath{\alpha}}
\newunicodechar{β}{\ensuremath{\beta}}
\newunicodechar{γ}{\ensuremath{\gamma}}
\newunicodechar{λ}{\ensuremath{\lambda}}
\newunicodechar{φ}{\ensuremath{\varphi}}
\newunicodechar{ε}{\ensuremath{\varepsilon}}
\newunicodechar{σ}{\ensuremath{\sigma}}
\newunicodechar{θ}{\ensuremath{\theta}}
\newunicodechar{δ}{\ensuremath{\delta}}
\newunicodechar{η}{\ensuremath{\eta}}
\newunicodechar{π}{\ensuremath{\pi}}
\newunicodechar{ξ}{\ensuremath{\xi}}
\newunicodechar{Σ}{\ensuremath{\Sigma}}
\newunicodechar{Δ}{\ensuremath{\Delta}}
\newunicodechar{∂}{\ensuremath{\partial}}
\newunicodechar{∀}{\ensuremath{\forall}}
\newunicodechar{∃}{\ensuremath{\exists}}
\newunicodechar{√}{\ensuremath{\surd}}
\newunicodechar{−}{\ensuremath{-}}
\newunicodechar{μ}{\ensuremath{\mu}}
\newunicodechar{ρ}{\ensuremath{\rho}}
\newunicodechar{τ}{\ensuremath{\tau}}
\newunicodechar{ω}{\ensuremath{\omega}}
\newunicodechar{Ω}{\ensuremath{\Omega}}
\newunicodechar{χ}{\ensuremath{\chi}}
\newunicodechar{ψ}{\ensuremath{\psi}}
\newunicodechar{Γ}{\ensuremath{\Gamma}}
\newunicodechar{Φ}{\ensuremath{\Phi}}
\newunicodechar{Λ}{\ensuremath{\Lambda}}
\newunicodechar{ℝ}{\ensuremath{\mathbb{R}}}
\newunicodechar{ℤ}{\ensuremath{\mathbb{Z}}}
\newunicodechar{ℕ}{\ensuremath{\mathbb{N}}}
\newunicodechar{△}{\ensuremath{\triangle}}
\newunicodechar{⊗}{\ensuremath{\otimes}}
\begin{document}
\section{2018年杭州电子科技大学601数学分析考研真题}\label{ux5e74ux676dux5ddeux7535ux5b50ux79d1ux6280ux5927ux5b66601ux6570ux5b66ux5206ux6790ux8003ux7814ux771fux9898}

来源：\texttt{HDU/601数学分析/2018/2018年杭州电子科技大学601数学分析考研真题.pdf}。原
PDF 页序：1, 2。

\section{杭州电子科技大学}\label{ux676dux5ddeux7535ux5b50ux79d1ux6280ux5927ux5b66}

\section{2018年攻读硕士学位研究生招生考试}\label{ux5e74ux653bux8bfbux7855ux58ebux5b66ux4f4dux7814ux7a76ux751fux62dbux751fux8003ux8bd5}

\section{《数学分析》试题}\label{ux6570ux5b66ux5206ux6790ux8bd5ux9898}

（试题共四大题，共2页，总分150分）

姓名 报考专业 准考证号

【所有答案必须写在答题纸上，做在试卷或草稿纸上无效！】

\subsection{一、填空题（本大题共8小题，每小题5分，本大题共40分）}\label{ux4e00ux586bux7a7aux9898ux672cux5927ux9898ux51718ux5c0fux9898ux6bcfux5c0fux98985ux5206ux672cux5927ux9898ux517140ux5206}

\begin{enumerate}
\def\labelenumi{\arabic{enumi}.}
\tightlist
\item
  \(\lim _ { n  \infty } ( 2 ^ { n } + 3 ^ { n } + 4 ^ { n } ) ^ { \frac { 1 } { n } } = (\)
  ）.
\end{enumerate}

2.设函数f(x)在x=1处可导，f(1)=3,则
\(\lim _ { x \to 0 } { \frac { f ( 1 + x ) - f ( 1 - 3 x ) } { x } } = (\)
）.

3.已知 \(e ^ { x + y } = \sin ( x y )\) 确定了隐函数y=y(x)，则
\({ \frac { d y } { d x } } = { }\) （ ）.

4.设 \(F ( x ) = \int _ { 0 } ^ { 2 x } x e ^ { t ^ { 2 } } d t\)
，则F'(x)=（ ）.

5.设 \(u = x ^ { 2 } + y - z ^ { 2 }\)
，则u在(1,-1,1)处的方向导数的最大值为（ ）

6．已知(x+ay)dx+(2x+y)dy为某函数的全微分，则a的值是（ ）

7．设函数f(x,y)连续，则二次积分
\(\int _ { - 2 } ^ { 0 } d x \int _ { 2 x + 4 } ^ { 4 - x ^ { 2 } } f ( x , y ) d y\)
可交换积分次序为（ ）

8.设M \(: x ^ { 2 } + y ^ { 2 } + z ^ { 2 } = 1\) ，则曲面积分
\(\iiint _ { \Sigma } ( 2 - x ^ { 2 } - y ^ { 2 } - z ^ { 2 } ) d S = (\)

\subsection{二、计算题（本大题共5小题，每小题6分，本大题共30分）}\label{ux4e8cux8ba1ux7b97ux9898ux672cux5927ux9898ux51715ux5c0fux9898ux6bcfux5c0fux98986ux5206ux672cux5927ux9898ux517130ux5206}

1．求极限值
\(\lim _ { x \to \infty } { \frac { x ^ { 3 } ( 1 - \cos { \frac { 1 } { x ^ { 2 } } } ) } { \sqrt { x ^ { 2 } + 1 } - x } }\)

2.设 \(u = u ( x , y ) , \nu = \nu ( x , y )\) 是由
\(\left\{ \begin{array} { l l } { \sin ( u + v ) = x y } \\ { e ^ { w } = 1 - x - y } \end{array} \right.\)
所确定的隐函数组，求
\(\frac { \hat { \alpha } u } { \hat { \alpha } x }\)

3.已知连续函数f(x)满足条件
\(\int _ { 0 } ^ { x } t f ( x - t ) d t = 1 - \cos x\) ，求
\(\int _ { 0 } ^ { \frac { \pi } { 2 } } f ( x ) d x\)

4.求
\(\scriptstyle \sum _ { n = 0 } ^ { \infty } { \frac { 1 } { n + 1 } } x ^ { n }\)
的和函数.

5.计算 \(\iiint z ^ { 2 } d \nu\) ，其中是由曲面
\(z = x ^ { 2 } + y ^ { 2 }\) 与 \(z = \sqrt { x ^ { 2 } + y ^ { 2 } }\)
围成的有界区域.

\subsection{三、综合题（本大题共3题，每题10分，本大题共30分）}\label{ux4e09ux7efcux5408ux9898ux672cux5927ux9898ux51713ux9898ux6bcfux989810ux5206ux672cux5927ux9898ux517130ux5206}

1.设f(0)=0，f''(0)存在，
\(g ( x ) = \left\{ { \begin{array} { l l } { { f ( x ) } } & { { x \not = 0 } } \\ { { x } } & { { } } \\ { { f ^ { \prime } ( 0 ) , } } & { { x = 0 } } \end{array} } \right.\)
.求 \(g ^ { \prime } ( 0 )\)

2．讨论函数项级数
\(\sum _ { n = 1 } ^ { \infty } ( - 1 ) ^ { n } { \frac { \sin ^ { 2 } x + n } { n ^ { 2 } } }\)
在 \(( - \infty , + \infty )\) 上的一致收敛性和绝对收敛性.

3．计算面积分
\(\iint _ { \Sigma } { x ^ { 2 } d y d z } + { y ^ { 2 } d x d z } + z ( 1 - 2 x - 2 y ) d x d y\)
，其中£为上半球面\(x ^ { 2 } + y ^ { 2 } + z ^ { 2 } = 1\) \(z \geq 0\)
的上侧.

\subsection{四、证明题（本大题共5题，每题10分，本大题共50分）}\label{ux56dbux8bc1ux660eux9898ux672cux5927ux9898ux51715ux9898ux6bcfux989810ux5206ux672cux5927ux9898ux517150ux5206}

1.设函数f(x)在(0,+∞)上满足f(x)=f(3x)且
\(\lim _ { x \to \infty } f ( x ) = 5\) ．证明在 \(( 0 , + \infty )\)
上\(f ( x ) \equiv 5\)

2.设
\(f ( x , y ) = \left\{ \begin{array} { c c } { { x ( x ^ { 2 } + y ^ { 2 } ) } } & { { ( x , y ) \neq ( 0 , 0 ) } } \\ { { x ^ { 2 } + y ^ { 4 } } } & { { ( x , y ) = ( 0 , 0 ) } } \end{array} \right.\)
，证明：

(1） \(\lim _ { ( x , y ) \to ( 0 , 0 ) } f ( x , y )\)
不存在；（2）f(0,0)和 \(f _ { y } ^ { \prime } ( 0 , 0 )\) 存在.

3．设连续函数列 \(\{ f _ { n } ( x ) \}\)
在闭区间{[}a,b{]}上一致收敛于f(x)．若 \(x _ { n } \in [ a , b ]\)
\(( n = 1 , 2 , 3 , \cdots )\) 且有
\(\lim _ { n \to \infty } x _ { n } = x _ { 0 }\) ．证明
\(\lim _ { n \to \infty } f _ { n } ( x _ { n } ) = f ( x _ { 0 } )\)

4.设f(x,y）与 \(\varphi ( x , y )\) 均在 \(( x _ { 0 } , y _ { 0 } )\)
点某邻域内可微，且 \(\varphi _ { y } ^ { \prime } ( x , y ) \neq 0\)
，已知\(( x _ { 0 } , y _ { 0 } )\) 是f(x,y）在约束条件
\(\varphi ( x , y ) = 0\)
下的一个极值点．试证当\(f _ { x } ^ { \prime } ( x _ { 0 } , y _ { 0 } ) \neq 0\)
时必有 \(f _ { y } ^ { \prime } ( x _ { 0 } , y _ { 0 } ) \neq 0\)

5．证明函数
\(F ( x ) = \int _ { 0 } ^ { + \infty } \frac { \sin { x t } } { 1 + t ^ { 2 } } d t\)
在 \([ 0 , + \infty )\) 上连续，并且在(0,+∞)上有连续的导函数.
\end{document}
